% Stele: A Lightweight Formal Verification Framework for Mathematical Reasoning
% LaTeX source -- Version 1.0.0
% Build: latexmk -pdf stele-whitepaper.tex
%        pdflatex stele-whitepaper.tex  (twice for references)

\documentclass[11pt,a4paper]{article}

\usepackage[T1]{fontenc}
\usepackage[utf8]{inputenc}
\usepackage{lmodern}
\usepackage[margin=2.5cm]{geometry}
\usepackage{amsmath,amssymb,amsthm}
\usepackage{microtype}
\usepackage{hyperref}
\usepackage{booktabs}
\usepackage{listings}
\usepackage{xcolor}
\usepackage{float}
\usepackage{enumitem}

% ---- Code listing style ----
\definecolor{codebg}{rgb}{0.97,0.97,0.97}
\lstdefinestyle{stele}{
  backgroundcolor=\color{codebg},
  basicstyle=\ttfamily\small,
  breaklines=true,
  frame=single,
  rulecolor=\color{gray!50},
  xleftmargin=8pt,
  xrightmargin=4pt,
  aboveskip=6pt,
  belowskip=6pt,
}
\lstset{style=stele}

% ---- Theorem environments ----
\theoremstyle{definition}
\newtheorem{definition}{Definition}
\newtheorem{remark}{Remark}
\theoremstyle{plain}
\newtheorem{claim}{Claim}

% ---- Macros ----
\newcommand{\stele}{\textsc{Stele}}
\newcommand{\kernel}{\texttt{kernel.py}}
\newcommand{\untrusted}{\textsc{untrusted}}
\newcommand{\todo}[1]{\textcolor{orange}{[TODO: #1]}}

\hypersetup{
  pdftitle={Stele: A Lightweight Formal Verification Framework},
  pdfauthor={Stele Project},
  colorlinks=true,
  linkcolor=blue!70!black,
  citecolor=green!40!black,
  urlcolor=blue!70!black,
}

% ====================================================================
\begin{document}
% ====================================================================

\title{\stele: A Lightweight Formal Verification Framework \\
  for Mathematical Reasoning}
\author{Stele Research Project \\[4pt]
  \small Version 1.0.0 \\
  \small \url{https://github.com/re-glow/stele-logic}}
\date{}
\maketitle

% ====================================================================
\begin{abstract}
% ====================================================================
We present \stele, a lightweight formal verification framework for
mathematical reasoning, designed as a small, auditable proof-checking
platform with explicit trust boundaries.
The core contribution is \emph{Stele-Light}, a natural-deduction
proof language in which users supply every inference step and the
trusted kernel verifies each step against a declared rule schema.
The trusted kernel is a 142-line Python module with zero runtime
dependencies.

Beyond the kernel, \stele{} provides: a Curry--Howard proof-term
core with bidirectional type checking and $\beta$-reduction for
intuitionistic propositional logic; structural diagnostics that
locate missing hypotheses, invalid transitions, undefined symbols,
and unused assumptions; matrix-valued semantic diagnostics for
Boolean, K3 (strong Kleene), and LP (Logic of Paradox) logics;
a bounded Kripke countermodel search for intuitionistic non-theorems;
versioned proof certificates with an independent re-verification
path; a browser-local Studio via Pyodide/WebAssembly; and an
optional statistical ML baseline for proof classification.
The system has zero runtime dependencies for the core verification
path and 1808 tests.

\stele{} is not a theorem prover, SMT/SAT solver, AI-powered
verifier, or full proof assistant.
All metatheoretic claims about the proof-term layer are supported
by regression and property-based tests, not by machine-checked
formal proofs.
\end{abstract}

% ====================================================================
\section{Introduction}
\label{sec:intro}
% ====================================================================

Students and researchers working with formal arguments --- in logic
courses, foundations of mathematics, or proof-theoretic investigations
--- need tools that are small, auditable, and honest about what they
verify.
Existing options occupy two extremes.
Full proof assistants such as Lean 4~\cite{Moura2021}, Rocq
(formerly Coq)~\cite{Rocq}, and Isabelle/HOL~\cite{Nipkow2002} are
powerful but require substantial learning investment, have large
codebases, and can be difficult to deploy or audit in educational
settings.
At the other extreme, AI-based proof assistance can generate
plausible-looking derivations without making the trust boundary
precise.

\stele{} occupies a different niche: a \emph{small, inspectable
proof-checking framework} with explicit, documented trust
boundaries.
The user writes every inference step; the kernel decides whether
each step is a valid instance of the declared rule.
The kernel does not search for proofs and does not use semantic
knowledge of the target logic --- it only matches step conclusions
against rule schemata syntactically.
This design makes the trust boundary precise and auditable:
\kernel{} is 142 lines of pure Python with no external dependencies.

Around this kernel, \stele{} provides several optional layers that
assist reasoning without entering the trust boundary: structural
diagnostics, a proof-term core for Curry--Howard exploration,
matrix semantic diagnostics, Kripke model search, and a
browser-local interactive Studio.
These layers are systematically labeled \untrusted{} and require
re-verification by the kernel for any claim of proof validity.

\subsection*{Contributions}

\begin{enumerate}[leftmargin=*]
  \item A small, dependency-free trusted proof kernel with a
        documented trust boundary and rule-data separation.
  \item A Curry--Howard proof-term calculus for intuitionistic
        propositional logic with bidirectional type checking,
        $\beta$-reduction, and a de Bruijn binder layer.
  \item A structural diagnostics platform with nine stable
        diagnostic codes and per-code explanations.
  \item Matrix semantic diagnostics for Boolean, K3, and LP
        logics, with per-rule soundness reports.
  \item A bounded Kripke countermodel search for propositional
        intuitionistic logic.
  \item Versioned proof certificates and an independent
        re-verification path (\texttt{minicheck}).
  \item A browser-local deployment via Pyodide/WebAssembly with
        no server-side computation.
  \item An optional ML baseline with corpus discipline,
        benchmark card, and measured metrics.
\end{enumerate}

% ====================================================================
\section{Positioning and Non-goals}
\label{sec:positioning}
% ====================================================================

\stele{} is a \textbf{proof checker and semantic diagnostics
platform}.
It is explicitly not:

\begin{itemize}
  \item \textbf{A theorem prover.}
        \stele{} does not search for proofs.
  \item \textbf{An SMT or SAT solver.}
        There is no constraint solving, satisfiability checking,
        or arithmetic reasoning.
  \item \textbf{An AI-powered verifier.}
        ML components are optional, isolated, and \untrusted{}.
  \item \textbf{A production proof assistant.}
        \stele{} lacks dependent types, a module system,
        higher-order unification, tactics, and theorem libraries.
  \item \textbf{A full first-order prover.}
        The proof-script surface is propositional.
  \item \textbf{A formally verified system.}
        Metatheoretic properties are supported by tests and
        proof sketches, not by machine-checked proofs.
\end{itemize}

\stele{} is:

\begin{itemize}
  \item A \textbf{rule-checked natural-deduction proof language}
        with a small, auditable trusted kernel.
  \item A \textbf{proof-term and metatheory playground} for
        exploring Curry--Howard correspondences.
  \item A \textbf{semantic diagnostics platform} for Boolean,
        K3, and LP matrix logics and bounded Kripke models.
  \item A \textbf{reproducible research software artifact} with
        zero runtime dependencies and a comprehensive test suite.
\end{itemize}

% ====================================================================
\section{Architecture and Trust Boundary}
\label{sec:arch}
% ====================================================================

\subsection{Design Principles}

The architecture follows two principles:

\begin{enumerate}
  \item \textbf{De Bruijn criterion.} The trusted code should be
        small enough for a human auditor to read completely.
        \kernel{} is 142 lines; the full trust surface including
        supporting modules (\texttt{ast.py}, \texttt{proof.py},
        \texttt{logic.py}, \texttt{parser.py}) is under 700 lines.

  \item \textbf{Rule-data separation.} Inference rules are not
        compiled into the kernel.
        They live in \texttt{logic.py} as \texttt{RuleSchema}
        frozen dataclass instances.
        Adding a new logic means adding a \texttt{RuleSchema}
        set; it never requires modifying the kernel.
\end{enumerate}

\subsection{Layer Diagram}

\begin{verbatim}
TRUSTED: kernel.py (142 lines, stdlib only)
  ast.py · proof.py · parser.py · logic.py · errors.py
  ------------------------------------------------
UNTRUSTED / OPTIONAL layers
  diagnostics.py   structural diagnostics
  proofgraph.py    dependency graph
  proofstate.py    proof-state hints
  matrix.py        many-valued semantics
  world.py         world/lattice demo
  kripke.py        Kripke countermodel search
  certificate.py   certificate emission
  minicheck.py     independent re-checker
  elaborate.py     script-to-term elaboration
  core/            proof-term calculus (IPL+IQL)
  browser.py       Pyodide bridge
  web.py           HTTP Studio server
  ------------------------------------------------
EXTERNAL / ISOLATED
  stele_ml/        optional ML baseline
  stele_lean/      optional Lean 4 bridge
\end{verbatim}

\subsection{The Kernel}

The kernel (\texttt{stele/kernel.py}) contains four functions:
\texttt{match}, \texttt{instantiate}, \texttt{check\_theorem},
and the internal \texttt{\_check\_block}.
\texttt{match} performs syntactic first-order pattern matching:
given a rule schema pattern and a concrete formula, it either
returns a substitution or \texttt{None}.
No semantic interpretation of connectives is used.
\texttt{check\_theorem} walks the proof tree, dispatching each
\texttt{have} step to a rule lookup and \texttt{match} call,
and each \texttt{suppose} block to a recursive check with a
scoped environment.

\subsection{Separation of Concerns}

A key invariant enforced by regression tests is that
\kernel{} does not import \texttt{matrix.py} and
\texttt{matrix.py} does not import \kernel{}.
Matrix semantics and proof-rule checking are strictly
separate concerns: a formula can be a tautology in one
matrix but not in another, independently of whether a
proof script demonstrates it.

% ====================================================================
\section{Stele-Light Proof Language}
\label{sec:language}
% ====================================================================

\subsection{Syntax}

A Stele-Light proof has the following structure:

\begin{lstlisting}
theorem <name> [using <logic>]:
  assume <label>: <formula>     -- introduce a hypothesis
  suppose <label>: <formula>    -- begin a dischargeable sub-proof
    have <label>: <formula> by <rule> <ref>...
    ...
  have <label>: <formula> by <rule> <ref>...
  conclude <formula> by <ref>
\end{lstlisting}

\noindent
\texttt{suppose} blocks introduce dischargeable assumptions.
When a \texttt{suppose} block closes, the assumption is no
longer available; the enclosing rule (e.g., \texttt{imp\_intro})
must reference both the suppose label and the block's conclusion.

\subsection{Formula Grammar}

\begin{align*}
  \text{formula} &\;::=\; \text{implication} \\
  \text{implication} &\;::=\; \text{disjunction}
    \;(\texttt{->}\;\text{formula})? \\
  \text{disjunction} &\;::=\;
    \text{conjunction}\;(\texttt{or}\;\text{conjunction})^* \\
  \text{conjunction} &\;::=\;
    \text{negation}\;(\texttt{and}\;\text{negation})^* \\
  \text{negation} &\;::=\; \texttt{not}\;\text{negation}
    \;\mid\; \text{atom} \\
  \text{atom} &\;::=\; \text{VAR} \;\mid\; \texttt{false}
    \;\mid\; \texttt{(}\;\text{formula}\;\texttt{)}
\end{align*}

Negation is sugar: $\lnot A \equiv A \to \bot$.
The parser rewrites \texttt{not}\;$A$ to
$\texttt{Op}(\texttt{"not"}, (A,))$; the type checker
normalizes this to $A \to \bot$ for type-level operations.

\subsection{Built-in Logics and Rules}

\begin{table}[H]
\centering
\caption{Built-in logics}
\begin{tabular}{@{}lll@{}}
\toprule
Logic & Mode & Notes \\
\midrule
\texttt{intuitionistic\_prop} & proof & 12 rules (IPL NJ) \\
\texttt{classical\_prop} & proof & 12 + 3: \texttt{dne}, \texttt{lem}, \texttt{pbc} \\
\texttt{K3} & matrix & Strong Kleene three-valued \\
\texttt{LP} & matrix & Logic of Paradox (Priest) \\
\texttt{boolean} & matrix & Classical Boolean \\
\bottomrule
\end{tabular}
\end{table}

The intuitionistic rule set comprises: \texttt{copy}, \texttt{mp},
\texttt{and\_intro}, \texttt{and\_elim\_left},
\texttt{and\_elim\_right}, \texttt{neg\_elim},
\texttt{ex\_falso}, \texttt{or\_intro\_left},
\texttt{or\_intro\_right}, \texttt{imp\_intro} (discharge),
\texttt{neg\_intro} (discharge), \texttt{or\_elim}
(two discharges).

\subsection{Example: Double-Negation Elimination}

\begin{lstlisting}
theorem dne_consequent:
  assume h1: not not P
  have h2: P by dne h1
  conclude P by h2
\end{lstlisting}

Accepted under \texttt{classical\_prop}; rejected under
\texttt{intuitionistic\_prop} because \texttt{dne} is absent
from the intuitionistic rule set.

% ====================================================================
\section{Proof-Term Core and Elaboration}
\label{sec:proofterms}
% ====================================================================

\subsection{Curry--Howard Correspondence}

The proof-term core (\texttt{stele.core}) implements a simply
typed $\lambda$-calculus corresponding to intuitionistic
propositional logic (IPL), with an experimental extension to
a first-order logic (FOL) fragment.
The correspondence is standard:

\begin{align*}
  A \to B &\;\longleftrightarrow\; \lambda x{:}A.\;e : A \to B \\
  A \land B &\;\longleftrightarrow\; \langle e_1, e_2 \rangle : A \times B \\
  A \lor B  &\;\longleftrightarrow\; \mathrm{inl}(e) \mathbin{/} \mathrm{inr}(e) : A + B \\
  \bot &\;\longleftrightarrow\; \text{empty type} \\
  \forall x.\,A(x) &\;\longleftrightarrow\; \Lambda x.\,e \text{ (experimental)}
\end{align*}

\subsection{Bidirectional Type Checking}

The type checker operates in two modes: a \emph{checking} mode
(given a type, verify the term has it) and an \emph{inference}
mode (derive the type of a term).
All constructors are frozen dataclasses; the checker traverses
them recursively.

\subsection{de Bruijn Binders}

Proof-variable binders use de Bruijn indices
(\texttt{stele.core.debruijn}).
The named API is preserved for user-facing code; the internal
representation uses indices to eliminate $\alpha$-renaming
issues in substitution.
Object-level variables in the FOL fragment remain named;
de Bruijn conversion at the object level is partially
implemented.

\subsection{Script-to-Term Elaboration}

\texttt{crosscheck\_theorem} provides a three-stage pipeline:
(1) kernel check of the proof script,
(2) elaboration to a Curry--Howard proof term,
(3) type check of the elaborated term.
All three must succeed; any failure is reported distinctly.
The elaboration path currently supports the full intuitionistic
rule set; classical rules are not elaborated.

\subsection{Reduction and Normalization}

\texttt{normalize} reduces a closed proof term via
$\beta$-reduction (call-by-name).
A fuel counter prevents non-termination; the default is 1000
steps.
$\eta$-reduction is not implemented.

\subsection{Metatheory Status}

The following claims are supported by proof sketches in
\texttt{docs/metatheory.md} and by regression tests, but are
\textbf{not machine-checked}:

\begin{enumerate}
  \item \textbf{Subject reduction}: well-typed terms remain
        well-typed after one reduction step.
        Status: proof sketch + regression tests + optional
        property-based tests (Hypothesis).
  \item \textbf{Normalization}: every closed IPL term
        normalizes.
        Status: follows from strong normalization for STLC;
        fuel limit provides a practical bound.
  \item \textbf{Confluence}: local confluence holds for
        $\beta$-reduction; Church--Rosser follows by Newman's
        lemma.
        Status: proof sketch; complete proof not machine-checked.
  \item \textbf{Consistency}: no closed term of type $\bot$
        can be constructed.
        Status: verified for the rule set by construction;
        regression tests confirm specific absence of
        proof-of-false terms.
\end{enumerate}

\begin{remark}
  Python regression tests are evidence of implementation
  correctness, not mathematical proof.
  A machine-checked metatheory would require formalizing these
  claims in a proof assistant such as Lean, Rocq, or Agda.
  This is a documented future direction.
\end{remark}

% ====================================================================
\section{Structural Diagnostics and Proof-State Tooling}
\label{sec:diagnostics}
% ====================================================================

\subsection{Structural Diagnostics}

\texttt{stele.diagnostics} performs multi-pass structural
analysis of a proof script outside the kernel.
Diagnostics are labeled \untrusted{} and do not constitute
proof verification.
Nine stable codes are defined:

\begin{table}[H]
\centering
\caption{Stable diagnostic codes}
\begin{tabular}{@{}ll@{}}
\toprule
Code & Description \\
\midrule
\texttt{UndefinedSymbol} & Variable used before introduction \\
\texttt{MissingHypothesis} & Referenced label not in scope \\
\texttt{UnsupportedConclusion} & Conclusion does not follow from steps \\
\texttt{CircularDependency} & Step references itself transitively \\
\texttt{UnusedAssumption} & Assumed formula never referenced \\
\texttt{UndefinedDefinition} & Definition references undeclared name \\
\texttt{InvalidTransition} & Rule application structurally inconsistent \\
\texttt{TypeMismatch} & Rule argument has wrong formula type \\
\texttt{KripkeCountermodelFound} & Conclusion has a Kripke countermodel \\
\bottomrule
\end{tabular}
\end{table}

Each code has a \texttt{DiagnosticExplanation} providing
a short description, likely cause, fix suggestion, and
example.

\subsection{Dependency Graph}

\texttt{stele.proofgraph} builds a labeled directed
dependency graph.
Each node is a labeled step; edges represent logical
dependencies.
Graph analysis detects cycles, unused assumptions, and
isolated steps.
DOT output is supported for external visualization.

\subsection{Proof-State Inspection}

\texttt{stele.proofstate} provides \untrusted{} proof-state
snapshots and rule-applicability hints.
\texttt{proof\_state\_from\_text} returns a \texttt{ProofState}
recording available labels and formulas, scope depths,
open assumptions, discharged labels, and the current pending
goal.
\texttt{suggest\_rule\_hints} applies 10 structural
patterns to suggest candidate next steps.
All hints have \texttt{trusted=False}; they must be
kernel-rechecked before use.

% ====================================================================
\section{Semantic Diagnostics}
\label{sec:semantics}
% ====================================================================

\subsection{Matrix Semantics}

A \texttt{Matrix} consists of a set of truth values,
designated values, and truth tables for the connectives.
Three matrices are built in:

\begin{table}[H]
\centering
\caption{Built-in matrices}
\begin{tabular}{@{}llll@{}}
\toprule
Matrix & Values & Designated & Notes \\
\midrule
\texttt{boolean} & \{T, F\} & \{T\} & Classical \\
\texttt{K3} & \{T, I, F\} & \{T\} & Strong Kleene~\cite{Kleene1952} \\
\texttt{LP} & \{T, I, F\} & \{T, I\} & Logic of Paradox~\cite{Priest1979} \\
\bottomrule
\end{tabular}
\end{table}

The K3 and LP tables follow standard conventions, fixed by test:
$I \to F = I$, $F \to I = T$.

\texttt{rule\_soundness(schema, matrix)} checks whether each
rule schema preserves designated values in all assignments.
Classical LEM is unsound in K3 (counterexample: $A{=}I$);
explosion is sound in K3 but not in LP.

\subsection{World and Lattice Demo}

\texttt{stele.world} is a toy propositional world system where
each \texttt{World} carries a matrix and an axiom set.
This is a pedagogical demonstration, not set-theoretic forcing
or a completeness theorem.

\subsection{Kripke Countermodel Search}

\texttt{stele.kripke} implements finite bounded Kripke model
search for intuitionistic propositional logic.
Given a formula, the searcher enumerates Kripke frames with up
to $n$ worlds (default $n=4$), checking the intuitionistic
forcing condition at each world.

For \texttt{P or not P} (LEM), the search finds:

\begin{lstlisting}
worlds: [0, 1]   order: reflexive + {0 <= 1}
  world 0: {}       (neither P nor not-P forced)
  world 1: {P}      (P is forced)
\end{lstlisting}

This correctly shows LEM is not intuitionistically valid.

\begin{remark}[Limitations]
  The search is bounded and incomplete.
  It may fail to find a countermodel that exists beyond the
  frame bound.
  The search is propositional only.
\end{remark}

% ====================================================================
\section{Proof Certificates and Minicheck}
\label{sec:certs}
% ====================================================================

\subsection{Certificate Format}

After a successful kernel verification, \stele{} can emit a
versioned JSON certificate recording each proof step as a
structured AST node.
The certificate format is versioned (\texttt{"version": "1"})
and includes the theorem name, logic, conclusion, and all
intermediate steps with their rules and references.

\subsection{Minicheck}

\texttt{stele.minicheck} re-verifies a certificate without
importing \texttt{stele.kernel} or \texttt{stele.parser}.
It implements its own minimal rule-matching logic from first
principles, operating only on the JSON certificate.
This provides a partial independence guarantee: a tampered
certificate that passes minicheck without having passed the
kernel would indicate a minicheck bug.

Tests include tamper detection (modifying step formulas,
removing steps, changing rule names) and cross-validation.

\begin{remark}
  Minicheck is Python and is not formally verified.
  A stronger guarantee would require a checker in a different
  language (Rust, OCaml) or a proof assistant.
  This is a documented future direction.
\end{remark}

% ====================================================================
\section{Browser Deployment and Distribution}
\label{sec:deployment}
% ====================================================================

\stele{} Studio runs the full trusted kernel inside the browser
via Pyodide~\cite{Pyodide} (Python compiled to WebAssembly).
No proof text is sent to any external server.
The Studio provides five interactive panels: Verify, Diagnose,
Graph, Semantic, and Proof State.

Four distribution modes are available:

\begin{description}
  \item[Hosted site] GitHub Pages, auto-deployed on push to \texttt{main}.
  \item[Single-file HTML] \texttt{stele.html} embeds the Stele
    source; Pyodide loaded from CDN.
  \item[Standalone executable] Built with PyInstaller (${\sim}$50\,MB).
  \item[Local Python] \texttt{python -m stele}; Python 3.10+,
    no extra dependencies.
\end{description}

The first-load Pyodide runtime is approximately 8\,MB, cached
by the browser after the initial visit.

% ====================================================================
\section{Optional ML Baseline and Corpus Discipline}
\label{sec:ml}
% ====================================================================

\subsection{Architecture}

\texttt{stele\_ml/} is an optional, isolated package providing
a statistical proof classification baseline.
It is entirely outside the trusted verification path
(enforced by \texttt{test\_ml\_isolation.py}).
The model is a Multinomial Naive Bayes classifier (Laplace
smoothing $\alpha{=}1.0$, bag-of-words features, pure Python
stdlib).

\subsection{Measured Metrics}

The following metrics are from the pipeline evaluation on the
committed 40-example synthetic sample (model trained on 400
in-memory examples, seed=0):

\begin{table}[H]
\centering
\caption{ML baseline metrics (40-example synthetic sample)}
\begin{tabular}{@{}lrl@{}}
\toprule
Metric & Value & Note \\
\midrule
Validity accuracy & 0.85 & 34/40 examples correct \\
Exact match & 0.60 & Validity AND all codes correct \\
Micro F1 (codes) & 0.50 & Weighted across predictions \\
Macro F1 (codes) & 0.36 & Averaged over codes with support \\
\bottomrule
\end{tabular}
\end{table}

\emph{These numbers must not be cited as general accuracy
claims.}
They are measurements on a 40-example synthetic corpus that
is not representative of real user proofs.

\subsection{Corpus Discipline}

Data discipline features include: manifest schema with
\texttt{label\_stats} and \texttt{creation\_command}; a
3-way train/dev/test split builder (\texttt{build\_dataset.py});
failure-mode analysis per code in evaluation reports; and a
benchmark card at \texttt{docs/benchmark-card.md} documenting
known limitations and reproduction steps.

% ====================================================================
\section{Evaluation and Testing}
\label{sec:testing}
% ====================================================================

\stele{} has 1808 tests (1804 passing, 4 skipped without
Hypothesis), organized across 44 test files.
Tests are run in CI on Python 3.10, 3.11, and 3.12.

Key test categories:

\begin{itemize}
  \item Kernel correctness:
    \texttt{test\_kernel\_valid.py} (correct proofs accepted),
    \texttt{test\_kernel\_invalid.py} (incorrect proofs rejected).
  \item Gallery honesty: \texttt{test\_gallery.py} verifies all
    15 curated examples against the live kernel on every CI run.
  \item Trust boundary: \texttt{test\_regression\_invariants.py}
    enforces mutual non-import between kernel/matrix and ML isolation.
  \item Certificates: \texttt{test\_minicheck.py} includes tamper
    detection and cross-validation tests.
  \item ML discipline: \texttt{test\_ml\_corpus\_discipline.py}
    checks manifest schema, split disjointness, report schema, and
    absence of overclaims.
  \item Property-based tests: \texttt{test\_proof\_term\_properties.py}
    (optional; requires Hypothesis).
\end{itemize}

The core CI workflow (\texttt{.github/workflows/ci.yml}) requires
only \texttt{pytest} --- no ML, LaTeX, or Lean dependencies.

% ====================================================================
\section{Related Work}
\label{sec:related}
% ====================================================================

\paragraph{Proof assistants.}
Lean 4~\cite{Moura2021}, Rocq/Coq~\cite{Rocq},
Isabelle/HOL~\cite{Nipkow2002}, and Agda~\cite{Norell2009}
are full proof assistants offering dependent types, rich tactic
languages, large theorem libraries, and machine-checked
metatheory.
\stele{} is not a replacement for these tools; it is designed
for settings where a small, auditable checker is more
appropriate.

\paragraph{Curry--Howard correspondence and proof terms.}
The Curry--Howard correspondence between intuitionistic proofs
and typed programs is foundational to the proof-term
core~\cite{SorensenUrzyczyn2006}.
\stele{}'s bidirectional type-checking architecture is
standard~\cite{PierceT98}.

\paragraph{Kripke semantics.}
Kripke models for intuitionistic propositional logic provide the
semantic underpinning for the countermodel search.
The forcing semantics is standard; see, e.g., \cite{TroelstraVanDalen1988}.

\paragraph{Many-valued logics.}
LP (Logic of Paradox) was introduced by Priest as a
paraconsistent logic in which contradictions do not trivialize
the system~\cite{Priest1979}.
K3 (strong Kleene) provides three-valued semantics for
sentences with truth-value gaps~\cite{Kleene1952}.
Both are matrix logics in the sense of many-valued logic
\cite{TODO:Malinowski}.

\paragraph{Proof certificates.}
The idea of emitting independently verifiable proof certificates
is well-established~\cite{TODO:LF}.
\stele{}'s certificate format is a lightweight JSON encoding;
the minicheck re-checker provides an independent software check.

\paragraph{ML for theorem proving.}
Several systems apply machine learning to guide tactic selection
or proof generation in full proof assistants~\cite{TODO:LeanDojo}.
\stele{}'s ML baseline is a proof \emph{classification}
model, not a proof-generation system.
It is explicitly \untrusted{} and evaluated only on a small
synthetic corpus.

% ====================================================================
\section{Limitations}
\label{sec:limitations}
% ====================================================================

\begin{itemize}
  \item \textbf{Propositional surface.}
    The Stele-Light proof language is propositional.
    No quantifier syntax exists at the script level.
  \item \textbf{No equality or function symbols.}
    The formula language has no arithmetic, equality, or
    function terms.
  \item \textbf{No proof search.}
    If the user does not supply every inference step, \stele{}
    cannot supply them.
  \item \textbf{Elaboration intuitionistic only.}
    The script-to-term elaboration supports only intuitionistic
    rules; classical proofs cannot be elaborated.
  \item \textbf{FOL experimental.}
    The FOL proof-term fragment is experimental; there is no
    surface syntax and object-variable de Bruijn conversion is
    incomplete.
  \item \textbf{No dependent types.}
    The proof-term calculus is simply typed.
  \item \textbf{Metatheory not machine-checked.}
    Subject reduction, normalization, confluence, and consistency
    are supported by proof sketches and tests, not by
    machine-checked proofs.
  \item \textbf{Kripke search bounded.}
    The countermodel search is bounded (${\le}4$ worlds default)
    and is neither sound nor complete as a proof procedure.
  \item \textbf{ML small corpus.}
    The ML baseline trains on a 40--400 record synthetic corpus;
    metrics on this corpus are not general accuracy claims.
  \item \textbf{Minicheck not formally verified.}
    The re-checker is an independent Python implementation, not
    a mechanically proved checker.
  \item \textbf{Browser requires CDN.}
    The Pyodide runtime (${\sim}$8\,MB) is fetched from CDN;
    fully offline mode is future work.
\end{itemize}

% ====================================================================
\section{Future Work}
\label{sec:future}
% ====================================================================

\begin{itemize}
  \item First-order surface syntax: adding quantifiers to
    Stele-Light with corresponding kernel rules.
  \item Equality and function symbols in the formula language.
  \item Structural rule policies (weakening, contraction,
    exchange) for linear, relevant, and hyperintensional logics.
  \item Rust or OCaml minicheck for a stronger independence
    guarantee.
  \item Machine-checked metatheory: formalizing subject
    reduction and confluence in Lean or Rocq.
  \item Lean bridge maturation: bidirectional import/export
    for a larger propositional fragment.
  \item ML data scaling after corpus discipline validation.
  \item Fully offline browser mode by bundling Pyodide.
  \item Kripke semantics for first-order logic.
\end{itemize}

% ====================================================================
\section{Conclusion}
\label{sec:conclusion}
% ====================================================================

\stele{} is a lightweight, auditable formal verification
framework occupying a niche between informal mathematical
notation and full proof assistants.
Its primary contribution is a small trusted kernel --- 142 lines,
zero runtime dependencies --- that checks natural-deduction
proofs against declared rule schemata syntactically.

Around this kernel, \stele{} provides optional layers for
proof-term extraction, structural diagnostics, matrix semantic
analysis, bounded Kripke countermodel search, proof certificates,
and browser-local deployment.
Every layer outside the kernel is labeled \untrusted{} and
documented as such.
Metatheoretic claims are supported by tests and proof sketches,
not machine-checked proofs, and this limitation is consistently
documented.

\stele{} v1.0.0 comprises 1808 tests, 15 curated gallery
examples with CI-enforced honesty labels, and zero runtime
dependencies for the core verification path.

% ====================================================================
\appendix
\section{Component Status Table}
\label{app:status}
% ====================================================================

\begin{table}[H]
\centering
\caption{Component status summary}
\small
\begin{tabular}{@{}p{4.5cm}lp{3.5cm}p{3cm}@{}}
\toprule
Component & Status & Evidence & Key Limitation \\
\midrule
Proof checker (Stele-Light) & Stable &
  1808 tests; kernel\_valid/invalid &
  Propositional only \\
Intuitionistic rules (12) & Stable &
  All kernel tests & --- \\
Classical extensions (dne/lem/pbc) & Stable &
  test\_classical\_principles & Not intuitionistic \\
Structural diagnostics & Stable &
  test\_diagnostics & UNTRUSTED; not all modes \\
Dependency graph & Stable &
  test\_proofgraph & DOT; needs Graphviz \\
Matrix semantics (K3/LP/bool) & Stable &
  test\_matrix, test\_rule\_soundness &
  Propositional; discharge skipped \\
World/lattice demo & Stable (demo) &
  test\_world, test\_world\_lattice & Toy demo; not forcing \\
Proof-term core (IPL) & Stable &
  test\_proof\_terms, test\_elaboration & Intuitionistic only \\
Proof-term reduction & Stable &
  test\_reduction & Fuel-limited; no $\eta$ \\
de Bruijn binders & Stable &
  test\_debruijn & Proof vars only \\
Proof-state hints & Stable (UNTRUSTED) &
  test\_proofstate (78) & Hints need recheck \\
Kripke countermodel search & Experimental &
  test\_kripke & Bounded; propositional \\
Proof certificates & Stable &
  test\_certificate & JSON; not formally verified \\
Minicheck & Stable &
  test\_minicheck & Python; not formally verified \\
Classical proof-term bridge & Experimental &
  test\_classical\_experimental & Formula translation only \\
FOL proof-term fragment & Experimental &
  test\_fol & No surface; incomplete \\
Browser Studio (Pyodide) & Stable &
  test\_studio & CDN required \\
Single-file HTML & Stable &
  test\_single\_html & CDN for Pyodide \\
Standalone executable & Stable &
  test\_packaging & PyInstaller to build \\
ML baseline & Optional / Experimental &
  test\_ml\_baseline & 40-record sample; UNTRUSTED \\
Lean 4 bridge & Optional / Experimental &
  test\_lean\_bridge & Propositional; needs Lean 4 \\
Metatheory (subject reduction etc.) & Proof sketches + tests &
  docs/metatheory.md & Not machine-checked \\
\bottomrule
\end{tabular}
\end{table}

% ====================================================================
\bibliographystyle{plain}
\bibliography{references}
% ====================================================================

\end{document}
